%% 306_Native_Zeit_Energie_Reziprozitaet_En_ch.tex
%% FFGFT -- Doc. 306: Replacing the Heisenberg import by native time-energy reciprocity
%% Author: Johann Pascher
%% Date: July 2026
%% Language: English

\section*{Doc.~306 \quad Native time--energy reciprocity --- replacing a borrowed justification}

\begingroup\small\noindent
\textbf{Abstract.} Documents~025 and~063 justify the absence of a cosmic initial singularity by
Heisenberg's energy--time uncertainty relation. That justification is here \emph{replaced}, not
refuted: the conclusion (no singular beginning) stands, the derivation is replaced by a native
one. Three defects are named --- an inverted implication (a finite $\Delta t$ yields a
\emph{finite} bound, not $\Delta E\to\infty$), a misapplication (the $\Delta t$ of the relation is
an observable's characteristic change time, not a cosmic age), and a breach of style
(\enquote{irrefutably} is not a ledger term). The replacing chain is native: from $T\cdot m = 1$
(Doc.~003/010), hence $T\cdot E = 1$, it follows that $E\to\infty \Leftrightarrow T\to 0$; but
FFGFT time carries a minimal, non-vanishing increment $d_1 = 100\,\xi_1 = 1/75$, fixed by
$\xi = 4/30000$. $T$ is therefore bounded below and $E$ bounded above --- an infinite density is
structurally excluded. In addition $\partial T^4 = \emptyset$: the question of a
\enquote{beginning} presupposes a boundary that a compact torus does not have.
\par\endgroup

\medskip

\section{What is being replaced}

Docs.~025 and~063 both contain the passage that Heisenberg's uncertainty relation
$\Delta E\cdot\Delta t \ge \hbar/2 = 1/2$ (natural units) \enquote{irrefutably proves} a classical
Big-Bang singularity to be impossible, and infer:
\begin{itemize}
\item a temporal beginning would mean $\Delta t$ finite;
\item this leads to $\Delta E \to \infty$;
\item hence the universe must have existed eternally.
\end{itemize}
The \emph{conclusion} --- no singular beginning, instead a finite core with a minimal scale $L_0$
from $\xi$ --- is retained. The \emph{derivation} is replaced. Three reasons.

\section{Three defects of the borrowed justification}

\paragraph{(M1) Inverted implication.} The inference \enquote{$\Delta t$ finite $\Rightarrow$
$\Delta E \to \infty$} is invalid. From $\Delta E \ge 1/(2\Delta t)$, a \emph{finite} $\Delta t$
yields a \emph{finite} lower bound on $\Delta E$; a beginning a finite time ago would give a tiny,
harmless bound. $\Delta E \to \infty$ follows only from $\Delta t \to 0$, i.e.\ from the
\emph{instant} $t=0$, not from a finite beginning. This is not a slip but a reversed implication.

\paragraph{(M2) Misapplication.} Energy--time uncertainty is not a canonical uncertainty relation.
By Pauli's theorem there is no self-adjoint time operator conjugate to a Hamiltonian bounded
below; the rigorous form (Mandelstam--Tamm) defines
$\Delta t = \Delta A / |\mathrm{d}\langle A\rangle/\mathrm{d}t|$ --- the characteristic
\emph{change time of an observable}, not an elapsed duration. The \enquote{age of the universe} is
not an admissible $\Delta t$ for this relation.

\paragraph{(M3) Breach of style.} \enquote{Irrefutably proves} is a strength of assertion that does
not occur in the FFGFT ledger. The ledger knows \emph{proven} (derived from $\xi$),
\emph{restricted}, and \emph{blocked}; a statement is certified as sufficiently precise to be
decidable, not as irrefutable.

\paragraph{The overarching point.} Heisenberg's relation is not false within the QM framework, but
it is \emph{projected}: it belongs to the description level that FFGFT understands as a projection
from the deeper structure. An argument against the cosmic singularity that leans on it does not
leave the projection --- it circulates within it. This is exactly what Doc.~298 warns of: one does
not exit the projection by importing its theorems.

\section{The replacing native chain}

\paragraph{(1) Exact reciprocity.} The time--mass duality (Doc.~003, cf.\ Doc.~010) reads
\begin{equation}
T \cdot m = 1, \qquad\text{and with } E = m \text{ (natural units)}\quad
\boxed{\,T \cdot E = 1\,}
\end{equation}
This is an \emph{equation}, not an inequality --- a sharper statement than
$\Delta E\cdot\Delta t \ge 1/2$, and a native one.

\paragraph{(2) Time has a minimal increment.} FFGFT time is the cumulated defect of the unrolled
winding,
\begin{equation}
D(k) = \sum_{n\le k} 100\,\xi_n \;\to\; \ln\!\left(1+\tfrac{k}{75}\right),
\qquad \xi_{n+1} = \xi_n\,(1-100\,\xi_n),
\end{equation}
with the first step
\begin{equation}
d_1 = 100\,\xi_1 = 100 \cdot \tfrac{4}{30000} = \tfrac{1}{75}.
\end{equation}
To be precise: $D(k)$ \emph{grows} without bound as $k\to\infty$ (logarithmically). The
singularity-free statement is not \enquote{$D$ never diverges} but: \emph{$D(k)$ is finite at
every finite $k$, and there is no point at which $D$ or its increments diverge.} The first step is
finite and fixed by $\xi$; no time increment is smaller than $d_1$.

\paragraph{(3) The inference.} From $T\cdot E = 1$,
\begin{equation}
E \to \infty \quad\Longleftrightarrow\quad T \to 0 .
\end{equation}
But no state with $T = 0$ exists in FFGFT: time is bounded below by a finite increment fixed by
$\xi$. Hence $E$ is bounded above --- an infinite energy density is \emph{structurally} excluded,
not by a statement about measurement precision. The exact numerical bound depends on the
normalisation and is not asserted here as a number; the structural statement stands without it.

\paragraph{(4) The boundary question.} The torus is compact and boundaryless,
$\partial T^4 = \emptyset$ (Doc.~302; cf.\ Doc.~263, Doc.~296). The question \enquote{what was at
$t=0$?} presupposes a boundary at which a beginning could sit. A compact torus has none. The
question is therefore not wrongly \emph{answered} but wrongly \emph{posed}: it imports the
decompactified projection as ontology. What standard cosmology reads as a singular beginning is
the point at which the \emph{projection} ends --- not the world.

\section{An observation on the genesis}

It is notable that 025 and 063 already point to $\xi$ in their \emph{result}: the singularity is
said to be \enquote{replaced by a tiny but finite core with minimal length scale $L_0$ (from
$\xi$)}. The answer therefore already came from $\xi$; only the ladder climbed to reach it was
borrowed. Why the import was made although Doc.~003/010 was available is a \emph{historical
question about the genesis of the corpus}, not a systematic one. It is marked as such here and
left unanswered; any answer would be speculation.

\section{Status}

\emph{Correction / replacement.} The conclusion (no singular beginning; a finite core with a
minimal scale set by $\xi$) is unchanged. The justification in Docs.~025 and~063 is \emph{no
longer authoritative} for this point; it is not recorded as refuted, but as superseded by the
native chain.

\emph{Proven (internal to FFGFT).} $T\cdot m = 1$ (Doc.~003/010); $D(k)$ and $d_1 = 1/75$ from
$\xi$ and the recursion (Doc.~295); $\partial T^4 = \emptyset$ (Doc.~302). The inference
$E\to\infty \Leftrightarrow T\to 0$ follows algebraically from the reciprocity.

\emph{Open.} The exact numerical upper bound on $E$ requires a fixed normalisation and is not
asserted here. The connection $R_H \to \ell_1$ remains open (P20, cf.\ Doc.~267); the cosmological
sector as a whole remains the least determined area of the corpus.
